\subsection{大组联调与结果比较}
各组先完成各自的方法调试，再在共同赛道中检查前后车之间的配合。联调时需要同时记录领行机器人和后车，尤其关注转角遮挡、前车停车、车距变化以及终点顺序。不同时间完成的单车试验可以帮助分析问题，但评价整体效果时仍需要一次连续运行。
\begin{longtable}{p{25mm}p{38mm}p{36mm}p{37mm}}
\caption{各部分实验情况与后续测试}\\\toprule
项目 & 主要检查内容 & 已有实验材料 & 后续需要完善的内容\\\midrule\endfirsthead
\toprule 项目 & 主要检查内容 & 已有实验材料 & 后续需要完善的内容\\\midrule\endhead
人形搬运 & 抓取、抱持、左右转与放置 & 动作组、现场照片和视频 & 全程重复运行、掉落与接触统计\\
方式一 & 标签识别、车距和转向 & AR实验步骤与控制说明 & 距离误差、转角丢失与停车时延\\
方式二 & 地图定位、绕障和路径保持 & 导航配置、地图与程序检查 & 矮障碍覆盖、完整路线跟随\\
方式三 & 身份保持、深度和局部导航 & RGB-D感知与控制程序 & 实际样本、遮挡和相似目标测试\\
方式四 & 车距、转弯和顺序停车 & 三车联调及问题记录 & 专门障碍与目标丢失场景\\
全队配合 & 同轮任务中各车到达与停车 & 多车现场画面与各组记录 & 连续运行记录和重复试验统计\\\bottomrule
\end{longtable}

\subsection{共同发现的问题}
各组在同一赛道调试时遇到了相互关联的问题。TonyPi的转弯和步态会影响后车标签的可见性；围栏较近时，局部避障容易与跟随方向冲突；前车长时间停车时，后车可能把它当成近障。方式二还需要检查深度扫描能否覆盖约35\,cm高的障碍，这一问题不能只通过调整地图参数解决。

我们将这些现象放在一起讨论，再分别修改动作、相机位置、跟随距离和速度限制。方式四的联调没有发生目标丢失，因此还不能评价实际重捕获效果；没有专门布置障碍的试验，也只能说明当时通道中的运行情况。后续需要增加相应场景，并保留成功和失败的记录。

\subsection{报告整理与小组配合}
报告整理时，我们尽量沿用各组原来的实验目的、步骤、控制公式和结果分析，只统一了符号、单位、图表编号与人员分工。第1组整理人形和地图导航部分，第2组整理AR实验，第3组整理RGB-D方案，第10组整理三车队列。各组相互检查前后接口，并共同完成总体设计、场地说明和实验总结。

这次合作让我们看到，同一个现场现象可能由多个环节共同造成。后车转向异常，不一定只是增益问题，也可能与标签朝向、前车步态或扫描触发有关。把各组观察到的情况放在一起，比各自只调单一参数更容易找出原因。
